\part{The system}

\section{Architecture}
\label{sec:arch}

\begin{figure}[h]
\centering
\begin{tikzpicture}[font=\small,
  box/.style={draw,rounded corners=2pt,align=center,fill=blue!4,minimum height=11mm,text width=24mm},
  io/.style={draw,align=center,fill=gray!12,minimum height=9mm,text width=22mm},
  fl/.style={-{Stealth[length=2mm]},thick}, bk/.style={-{Stealth[length=2mm]},dashed}]
\node[io] (g) at (0,0) {typed grammar\\+ schema\\+ data profile};
\node[box] (s1) at (3.2,0) {\textbf{1} derive test case (trees, weights)};
\node[box] (s2) at (6.4,0) {\textbf{2} attach oracle checks (tree rewrites)};
\node[box] (s3) at (9.6,0) {\textbf{3} run: server, plan variants, 18.6};
\node[box] (s4) at (12.8,0) {\textbf{4} classify, triage, reduce, localize};
\node[io] (c) at (0,-2.6) {recent \pg commit};
\node[box] (s5) at (3.2,-2.6) {\textbf{5a} LLM: fragment, seeds, rewrites};
\node[box] (s6) at (6.4,-2.6) {\textbf{5b} validation gate on the server};
\draw[fl] (g)--(s1); \draw[fl] (s1)--(s2); \draw[fl] (s2)--(s3); \draw[fl] (s3)--(s4);
\draw[bk] (s3.north) -- ++(0,0.5) -| node[pos=0.25,above]{\footnotesize feedback: corpus, weights, plan shapes} (s1.north);
\draw[fl] (c)--(s5); \draw[fl] (s5)--(s6);
\draw[bk] (s6.south) -- ++(0,-0.45) -| node[pos=0.25,below]{\footnotesize repair round} (s5.south);
\draw[fl] (s6.north) -- ++(0,0.6) -| (s1.south);
\end{tikzpicture}
\caption{The fuzzing loop (top) and the LLM loop (bottom). Steps match Sections~\ref{sec:grammar}--\ref{sec:llm}.}
\end{figure}

\subsection{Repository map}

\begin{tabularx}{\linewidth}{lX}
\toprule
Path & Contents \\
\midrule
\file{src/grammar/postgres.pgg} & the typed grammar \\
\file{src/corpus/schema.sql, schema.yaml} & the 9-relation schema (SQL for the server, YAML for the generator) \\
\file{src/corpus/rls.sql, profile.json} & RLS policies; data profile \\
\file{src/pgfuzz/grammar/} & DSL parser (\file{dsl.py}), engine (\file{engine.py}), generator, rewriter (oracles), data profile, tree I/O \\
\file{src/pgfuzz/reference/} & \pg executor, cross-version executor \\
\file{src/pgfuzz/eval/} & oracle, differential digests and determinism filter, signatures, metrics, statistics, features, known-bug catalog \\
\file{src/pgfuzz/triage/} & tokenizer, token reduction, tree reduction, localization, triage driver, precision \\
\file{src/pgfuzz/llm/} & commit miner, prompts, client, gate, diff coverage, loop \\
\file{src/pgfuzz/bench/} & re-injected bug benchmark (verify, attribute) \\
\file{src/pgfuzz/runner/, runlog/, report/} & experiment runner, run logs, reports \\
\file{src/generator/} & the random/mutation arm (Stage 1) \\
\file{src/bench/reinjected/} & benchmark: 23 commits, PoCs, \file{bugs.yaml}, \file{verify.json} \\
\file{src/configs/} & \file{rq1.yaml}, \file{rq3.yaml}, hunt and local configs \\
\file{scripts/} & server builds and starts (\file{build_reinjected.sh}, \file{rq1_servers.sh}, \file{rq3_servers.sh}, \file{ref_server.sh}), bisection, setup \\
\file{triage/verdicts.yaml} & the verdict of every triaged candidate (feeds RQ2) \\
\file{docs/} & guides, plans, evidence, bug reports (\file{docs/bughunt/}), paper draft PDF \\
\file{paper/} & the ACM/FSE paper (LaTeX) \\
\bottomrule
\end{tabularx}

\subsection{One trial}

\texttt{python -m pgfuzz run configs/rq1.yaml -j 4} plans the trials (arm
$\times$ trial), and each worker runs one trial against its own server
(port = base port + worker). A trial loops: ask the generator for a test case,
execute it, classify the result, update metrics and the run log, save a
failure if the signature is new, give feedback to the generator. Per trial it
writes \file{manifest.json}, \file{events.jsonl}, \file{timeline.jsonl}
(snapshots every 30 s), \file{failures/<signature>/} and \file{summary.json}.
Trials are resumable: re-running the command skips finished trials. The seed of
trial $k$ is the same in every arm (paired design).

\section{Step 1: the typed grammar and the generator}
\label{sec:grammar}

\subsection{The \texttt{.pgg} format}

\begin{itemize}
  \item \texttt{\%type int "int" num ord key \dots{} w=6} declares a SQL type,
    its classes (\texttt{num}, \texttt{ord}, \texttt{key}, \dots) and a weight.
  \item A rule has alternatives with weights \texttt{<4>}; terminals are
    keywords or quoted strings; \texttt{x?0.3} is optional with probability
    0.3; \texttt{x\{1,4/","\}} repeats 1--4 times separated by commas.
  \item Rules can be typed: \texttt{expr<T>}, \texttt{expr<int>},
    \texttt{expr<X:key>} (a type variable restricted to a class).
  \item Attributes: \texttt{@scope} opens a new name scope,
    \texttt{@set(agg\_ok)} sets a context flag (aggregates allowed),
    \texttt{\^{}} marks elements derived early (FROM before SELECT, so the
    select list can see the FROM's columns).
  \item Builtins (Python): \texttt{@col(T)} a visible column of type T,
    \texttt{@lit(T)} a literal, \texttt{@dlit(T)} a literal from the data of the
    compared column, \texttt{@table} a relation with a fresh alias,
    \texttt{@items} the select list, \texttt{@group\_keys} makes the level
    grouped, \texttt{@sorted\_agg} order-sensitive aggregates made
    deterministic, and about 20 more for DDL and DML.
\end{itemize}

A real excerpt:
\begin{lstlisting}
query @scope
  : <8> select_core order_by?0.3 limit_clause?0.06
  | <2> with_clause select_core order_by?0.2
  | <2> set_query
  ...
select_core
  : <6> SELECT distinct?0.12 @items ^from_clause ^where_clause?0.6
  | <3> SELECT @items ^from_clause ^where_clause?0.5 GROUP BY ^@group_keys having_clause?0.4
  | <2> SELECT @items ^from_clause ^where_clause?0.5 ^@global_agg having_clause?0.2
  ...
expr<bool>
  : <4> "(" expr<X:key> cmp expr<X> ")"
  | <1.2> "(" @col(X:ord) cmp @dlit(X) ")"
  | <2> "(" expr<bool> AND expr<bool> ")"
  | <1> "(" expr<X:key> NOT?0.3 IN "(" expr<X>{1,4/","} ")" ")"
  ...
\end{lstlisting}

Conventions that make oracles possible: every binary operator is
parenthesised (precedence never changes types); every column is qualified with
a unique alias; order-sensitive aggregates order by their own argument.

\subsection{The engine}

The engine expands a rule into a tree, choosing an alternative by weight,
binding type variables (\texttt{X}) consistently, tracking scopes (which
aliases and columns are visible, correlated or isolated) and flags (grouped
level, inside an aggregate). It backtracks when an alternative cannot be
completed (\texttt{Fail}) and prefers shallow alternatives past a soft depth.
Each alternative knows its minimal derivation depth (\texttt{mind}). While
rendering, each node's tokens are wrapped in invisible markers
(\verb|\x01<nid>| \dots{} \verb|\x02<nid>|), so the exact token span of every node
is known afterwards. A \texttt{MINIMAL} marker on a guide node makes the
engine derive that subtree with its smallest alternatives (used by tree
reduction).

\paragraph{Replay.} Given a guide tree, the engine follows its choices
(same alternatives, same columns, same literals), re-binding names and types
in the current schema state. Mutation and reduction both use replay, so their
results stay valid. Trees are stored with failures as JSON using stable
alternative \emph{ordinals} (uids are process-local) and exact container
types; 599 of 600 test cases replay to identical SQL in a fresh process.

\subsection{The generator}

\paragraph{A test case} is 0--2 setup statements (DDL such as
\texttt{CREATE INDEX}/\texttt{TABLE}, DML, \texttt{ANALYZE}, \texttt{SET LOCAL}
of a planner setting) followed by 1--2 queries or DML statements.

\paragraph{Fresh or mutated.} With probability 0.5 a corpus entry is picked
(weighted by its finds and how often it was picked) and mutated:
\emph{regen} (re-derive one subtree), \emph{splice} (graft a subtree from
another corpus entry), \emph{insert} (add a statement), \emph{drop} (remove one).

\paragraph{Corpus admission.} A test case enters the corpus if it produced a
new bug signature, a new SQL feature (catalogue of 148), a new pair of
features, or a new plan shape (digest of node types of
\texttt{EXPLAIN (FORMAT JSON, COSTS OFF)}).

\paragraph{Feedback-weighted alternatives} (plan 5.3). Each alternative has a
weight multiplier. After every test case, each alternative used in a tree gets
reward $r = 1$ if the test case entered the corpus, $-0.5$ if that tree's
statement failed, else 0; its weight is multiplied by $e^{\eta(r - \bar r)}$
with $\eta=0.1$ and $\bar r$ a running average, clipped to $[0.2, 5]$.
\emph{Measured:} 533 vs 419 distinct plan shapes in 1{,}000 test cases at the
same validity; the bandit learned to damp window-function alternatives that
often fail. RQ1 contains an ablation arm without it.

\paragraph{Data profile} (plan 6, interval and constraint-based analysis).
\texttt{pgfuzz profile} asks the server for min, max and the 10/25/50/75/90\%
quantiles of every column (\file{corpus/profile.json}). With probability 0.3 a
literal comes from it: boundary values (observed min$\pm$1, max$\pm$1, 0, $-0$,
NaN, $\pm$Infinity, the largest int, the empty array) or quantiles. The new
builtin \texttt{@dlit(T)} takes a quantile of the \emph{column the predicate
compares} (the engine remembers the last \texttt{@col}), so
\texttt{col < @dlit} selects a real share of rows. \emph{Measured:}
column-vs-constant predicates selecting 5--95\% of their table: 29.3\%
$\rightarrow$ 58.6\%.

\paragraph{Validity tuning loop.} \texttt{pgfuzz sample --execute} prints the
valid rate, the top error messages and the worst alternatives (by valid
rate); each fix is to one rule. This loop raised validity from 83\% to 100\%
before meeting 1.

\subsection{The schema}

Nine relations shared by both arms: \texttt{t0} (int, text, float8), \texttt{t1}
(int PK, int, jsonb, int4range), \texttt{t2} (identity bigint, numeric,
timestamptz, text[]), \texttt{t3} (int NOT NULL, bool, date, interval, int[],
text), \texttt{t4} (virtual and stored generated columns), \texttt{tp}
(range-partitioned), \texttt{tl} (list-partitioned), \texttt{ti} (inheritance
parent with child), \texttt{v0} (a view). Data is deterministic (no
\texttt{now()}), with NULLs in most columns. Indexes: B-tree (including
multi-column \texttt{t3(c0,c2)} for skip scan), GIN, GiST, BRIN, and an
expression index with a predicate.
